\section{Tilgungsrechnung}

\stepcounter{subsection}% Im Original steht vor den Beispielen kein eigener Unterabschnitt (Nummerierung 1.3.1.x)

Die Rückzahlung von Krediten geschieht meist in einer Form, bei der der Kreditnehmer den Schuldbetrag durch regelmäßige Zahlungen ausgleicht. Sind diese Zahlungen alle gleich hoch (dies nennt man \textbf{Annuitätentilgung}), so befindet man sich finanzmathematisch wieder in der Rentenrechnung -- mit dem Unterschied, dass man die Beträge sozusagen aus »anderer Perspektive betrachtet«. Zum Teil werden auch andere Bezeichnungen verwendet.

\subsubsection{Beispiel}

Ein Darlehen in Höhe von 140.000\,€ wird mit monatlich 0{,}8\,\% verzinst -- dies entspricht einem nominellen Zinssatz von 9{,}6\,\%. Wie hoch ist die monatliche Annuität bei einer Rückzahlungsdauer von 10 Jahren?

Man kann die Frage so auffassen: Die anfängliche Schuld $S$ von 140.000\,€ entspricht dem Rentenbarwert, und durch die regelmäßigen Zahlungen der Raten $r$, die hier als Annuitäten $A$ bezeichnet werden, soll die in 10 Jahren aufgelaufene Schuld abgetragen werden. Dies führt mit der bekannten Formel für den Rentenendwert und mit den geänderten Bezeichnungen zu
\[ S \cdot q^n = A \cdot \frac{q^n - 1}{q - 1} \]
bzw. nach $A$ aufgelöst:
\[ A = S \cdot q^n \cdot \frac{q - 1}{q^n - 1} \]
Hier ergibt sich somit:
\[ A = 140.000 \cdot 1{,}008^{120} \cdot \frac{1{,}008 - 1}{1{,}008^{120} - 1} \approx 1.819{,}24 \]
Es sind also monatlich 1.819{,}24\,€ zu zahlen.\beispielende

In der Praxis -- beispielsweise bei der Aufnahme eines Baukredits -- wird meist nicht die Laufzeit, sondern die Annuität vorgegeben. Die Laufzeit kann dann berechnet werden. Die Annuität kann etwa davon abhängig gemacht werden, einen wie hohen Betrag der Schuldner monatlich aufbringen kann. Ein anderes gebräuchliches Verfahren setzt die Annuität so fest, dass z.\,B. im ersten Jahr ein gewisser Prozentsatz der Schuld getilgt wird.

\subsubsection{Beispiel}

Ein Baudarlehen über 250.000\,€ wird zu 6\,\% Jahreszins aufgenommen. Da im ersten Jahr 15.000\,€ Zinsen zu zahlen sind und 1\,\% der Darlehenssumme getilgt werden soll, wird eine Annuität von 17.500\,€ festgelegt, die konstant gezahlt werden soll. Nach welcher Zeit ist das Darlehen vollständig abbezahlt?

Die Grundformel
\[ S \cdot q^n = A \cdot \frac{q^n - 1}{q - 1} \]
kann nach $q^n$ aufgelöst werden, dies ergibt:
\[ q^n = \frac{A}{A - S(q - 1)} \qquad \text{(Bitte nachrechnen!)} \]
Bei genauerem Hinsehen erkennt man, dass die rechte Seite gerade das Verhältnis von Annuität zur Tilgung im ersten Jahr ist, denn $S \cdot (q - 1)$ sind gerade die für das erste Jahr zu zahlenden Zinsen. Einsetzen der Werte führt hier zu
\[ 1{,}06^n = 7 \]
und somit
\[ n = \frac{\lg 7}{\lg 1{,}06} \approx 33{,}4 . \]
Das Darlehen ist also nach 34 Jahren getilgt, wobei im letzten Jahr nicht mehr die volle Annuität anfällt.\beispielende

In der Tilgungsrechnung ist es auch üblich, einen sogenannten \textbf{Tilgungsplan} aufzustellen, in dem für die fortlaufenden Jahre jeweils die gezahlten Zinsen, die erfolgte Tilgung und die verbleibende Restschuld aufgeführt sind. Mit einem Tabellenkalkulationsprogramm wie Microsoft Excel lässt sich ein solcher Tilgungsplan leicht erstellen. Für das obige Beispiel ist der Tilgungsplan auf der nächsten Seite dargestellt. Wie zu sehen, ist im 34. Jahr nur noch eine Restschuld von 6.642{,}09\,€ verblieben.

\begin{center}
  \centering
  \textbf{Tilgungsplan zu Beispiel 1.3.1.2 (Beträge in €):}\par\medskip
  \begin{tabular}{rrrrr}
    \hline
    Jahr & Restschuld & Zinsen & Tilgung & Annuität\\
    \hline
    1 & 250.000{,}00 & 15.000{,}00 & 2.500{,}00 & 17.500\\
    2 & 247.500{,}00 & 14.850{,}00 & 2.650{,}00 & 17.500\\
    3 & 244.850{,}00 & 14.691{,}00 & 2.809{,}00 & 17.500\\
    4 & 242.041{,}00 & 14.522{,}46 & 2.977{,}54 & 17.500\\
    5 & 239.063{,}46 & 14.343{,}81 & 3.156{,}19 & 17.500\\
    6 & 235.907{,}27 & 14.154{,}44 & 3.345{,}56 & 17.500\\
    7 & 232.561{,}70 & 13.953{,}70 & 3.546{,}30 & 17.500\\
    8 & 229.015{,}41 & 13.740{,}92 & 3.759{,}08 & 17.500\\
    9 & 225.256{,}33 & 13.515{,}38 & 3.984{,}62 & 17.500\\
    10 & 221.271{,}71 & 13.276{,}30 & 4.223{,}70 & 17.500\\
    11 & 217.048{,}01 & 13.022{,}88 & 4.477{,}12 & 17.500\\
    12 & 212.570{,}89 & 12.754{,}25 & 4.745{,}75 & 17.500\\
    13 & 207.825{,}15 & 12.469{,}51 & 5.030{,}49 & 17.500\\
    14 & 202.794{,}66 & 12.167{,}68 & 5.332{,}32 & 17.500\\
    15 & 197.462{,}34 & 11.847{,}74 & 5.652{,}26 & 17.500\\
    16 & 191.810{,}08 & 11.508{,}60 & 5.991{,}40 & 17.500\\
    17 & 185.818{,}68 & 11.149{,}12 & 6.350{,}88 & 17.500\\
    18 & 179.467{,}80 & 10.768{,}07 & 6.731{,}93 & 17.500\\
    19 & 172.735{,}87 & 10.364{,}15 & 7.135{,}85 & 17.500\\
    20 & 165.600{,}02 & 9.936{,}00 & 7.564{,}00 & 17.500\\
    21 & 158.036{,}02 & 9.482{,}16 & 8.017{,}84 & 17.500\\
    22 & 150.018{,}18 & 9.001{,}09 & 8.498{,}91 & 17.500\\
    23 & 141.519{,}27 & 8.491{,}16 & 9.008{,}84 & 17.500\\
    24 & 132.510{,}43 & 7.950{,}63 & 9.549{,}37 & 17.500\\
    25 & 122.961{,}06 & 7.377{,}66 & 10.122{,}34 & 17.500\\
    26 & 112.838{,}72 & 6.770{,}32 & 10.729{,}68 & 17.500\\
    27 & 102.109{,}04 & 6.126{,}54 & 11.373{,}46 & 17.500\\
    28 & 90.735{,}59 & 5.444{,}14 & 12.055{,}86 & 17.500\\
    29 & 78.679{,}72 & 4.720{,}78 & 12.779{,}22 & 17.500\\
    30 & 65.900{,}50 & 3.954{,}03 & 13.545{,}97 & 17.500\\
    31 & 52.354{,}53 & 3.141{,}27 & 14.358{,}73 & 17.500\\
    32 & 37.995{,}81 & 2.279{,}75 & 15.220{,}25 & 17.500\\
    33 & 22.775{,}55 & 1.366{,}53 & 16.133{,}47 & 17.500\\
    34 & 6.642{,}09 & & & \\
    \hline
  \end{tabular}
\end{center}

\subsubsection{Übungsaufgabe}

\begin{graukasten}
  Ein Darlehen von 50.000\,€ soll bei einem Zinssatz von 8\,\% p.\,a. in 15 gleich hohen nachschüssigen Annuitäten getilgt werden. Wie viel muss jeweils am Jahresende gezahlt werden?
\end{graukasten}

Neben der Annuitätentilgung gibt es auch die Möglichkeit der \textbf{Ratentilgung}, bei der nicht die Höhe der jährlichen Zahlungen konstant ist, sondern der jährlich getilgte Betrag. Da die Restschuld jährlich kleiner wird und somit immer weniger Zinsen zu zahlen sind, wird auch der jährlich zu zahlende Betrag immer kleiner. Daher kann man auf diese Form der Tilgung nicht mehr die Formeln der Rentenrechnung anwenden. Selbstverständlich kann man jedoch auch in diesem Fall leicht einen Tilgungsplan aufstellen.

\subsubsection{Beispiel}

Jemand hat am 01.01.2011 einen Kredit über 60.000\,€ aufgenommen. Er verpflichtet sich, diese Schuld in 6 gleich großen am Jahresende zu zahlenden Tilgungsraten in Höhe von 10.000\,€ zu tilgen. Dazu kommen jeweils noch die aufgelaufenen Zinsen, wobei ein Zinssatz von 8\,\% angewendet wird. Weitere Gebühren gibt es nicht. Der folgende Tilgungsplan gibt einen Überblick über die anfallenden Zahlungen und die Restschuldentwicklung:

\begin{center}
  \centering
  \begin{tabular}{rrrrr}
    \hline
    Jahr & Restschuld (zu Beginn des Jahres) & Zinsen $Z_j$ & Tilgung $T_j$ & Annuität $A_j$\\
    \hline
    1 & 60.000 & 4.800 & 10.000 & 14.800\\
    2 & 50.000 & 4.000 & 10.000 & 14.000\\
    3 & 40.000 & 3.200 & 10.000 & 13.200\\
    4 & 30.000 & 2.400 & 10.000 & 12.400\\
    5 & 20.000 & 1.600 & 10.000 & 11.600\\
    6 & 10.000 & 800 & 10.000 & 10.800\\
    7 & 0 & & & \\
    \hline
  \end{tabular}
\end{center}
\beispielende

Selbstverständlich können auch bei der Ratentilgung Formeln angegeben werden, mit deren Hilfe Restschuld, Zinsen usw. im jeweiligen Jahr ausdrückbar sind -- darauf soll hier jedoch verzichtet werden.

\subsubsection{Übungsaufgabe}

\begin{graukasten}
  Ein Darlehen von 30.000\,€ soll jährlich mit 20\,\% (vom Anfangsdarlehen) getilgt werden. Der Zinssatz beträgt 7{,}5\,\%. Man erstelle einen Tilgungsplan.
\end{graukasten}
